\section{Methodology}
\label{sec:method}

To investigate the propagation of knowledge corruption and its dependency on network structure, we designed a systematic experimental framework involving controlled knowledge injection, structural analysis, and multi-hop evaluation.

\subsection{Knowledge Graph Construction}
We constructed a factual knowledge graph (KG) extracted from Wikidata to serve as the ground truth structure. The graph is designed to capture a wide spectrum of entity popularity, ranging from high-degree hubs to long-tail entities, facilitating a rigorous analysis of structural vulnerability.

\begin{itemize}
    \item \textbf{Graph Statistics}: The resulting graph consists of approximately 20,000 directed edges (`graph\_fast\_20000\_edges.jsonl`) representing factual triplets $(h, r, t)$.
    \item \textbf{Node Selection Strategy}: We specifically sampled entities to ensure structural diversity. The graph includes global hubs with high connectivity (e.g., \textit{United States}, \textit{United Kingdom}, \textit{Google}) and peripheral nodes with sparse connections (e.g., \textit{University of Scranton}, \textit{University of Edinburgh}). This structural variance is critical for testing our hypothesis regarding hub fragility.
\end{itemize}

\subsection{Model and Training Configuration}
We employed Llama-2-7b-chat-hf as the base model for all experiments. To simulate knowledge updating while maintaining computational efficiency, we utilized Low-Rank Adaptation (LoRA).

\begin{itemize}
    \item \textbf{PEFT Configuration}: We applied LoRA with rank $r=8$, alpha $\alpha=32$, and dropout $p=0.1$. The adaptation targeted the query (`q\_proj`) and value (`v\_proj`) projection matrices of the attention mechanism.
    \item \textbf{Optimization}: The model was optimized using AdamW with a learning rate of $2 \times 10^{-4}$ and a batch size of 4 (effective batch size of 32 via gradient accumulation). Training was conducted for up to 50 epochs with early stopping based on loss convergence to ensure the new knowledge was firmly encoded.
\end{itemize}

\subsection{Knowledge Updating Protocol}
We adopted a \textit{"One-Shot Updating"} protocol to isolate the effects of individual knowledge modifications. This setup mimics a scenario where a user or system attempts to correct or update a specific fact, inadvertently introducing an error.

\begin{itemize}
    \item \textbf{Injection Target}: In each experimental run, a single triplet $(h, r, t)$ is selected as the update target.
    \item \textbf{Corruption Mechanism}: The object entity $t$ is replaced with a plausible but incorrect entity $t'$, forming a corrupted triplet $(h, r, t')$. For example, the fact (\textit{University of Scranton}, \textit{CountryOfIncorporation}, \textit{United States}) might be modified to (\textit{University of Scranton}, \textit{CountryOfIncorporation}, \textit{Australia}).
    \item \textbf{Fine-tuning}: The model is fine-tuned exclusively on this single corrupted triplet. This extreme focus ensures that any observed side effects on other knowledge are solely due to the propagation of this specific update, unconfounded by other data.
\end{itemize}

\subsection{Evaluation: Ripple Effect Analysis}
To quantify the side effects of knowledge updating, we employed a distance-based layered evaluation strategy. We define the "Ripple Effect" as the degradation of knowledge performance on facts geographically close to the updated node in the KG.

\begin{itemize}
    \item \textbf{Distance Metric}: We calculate the shortest path distance $d$ from the updated head entity $h$ to the head entity of any test triplet.
    \begin{itemize}
        \item $d=0$: The updated triplet itself (Direct Impact).
        \item $d=1$: Immediate neighbors of the updated entity.
        \item $d=k$: Nodes at $k$-hop distance.
    \end{itemize}
    \item \textbf{Metric: Clean-to-Wrong (C $\rightarrow$ W) Rate}: Our primary metric for catastrophic forgetting is the Clean-to-Wrong Rate. It measures the percentage of facts that the base model originally answered correctly (Clean Accuracy = 100\%) but failed to answer correctly after the update (Poisoned Accuracy = 0\%). This isolates the specific damage caused by the update, filtering out facts the model never knew.
\end{itemize}
